\chapter{Introduzione: cos'è l'IA e cosa fa questo progetto}\label{cap:introduzione}

% Capitolo 1 — Introduzione. Cos'è l'IA e cosa fa il progetto supervised_telemedicina.

Questo documento spiega, dal basso, il progetto \progetto{}: un sistema di
intelligenza artificiale per la telemedicina che classifica sei parametri
vitali in tre classi di rischio. È pensato per chi sa programmare ma non ha
mai studiato l'IA: useremo spesso analogie con il codice che già conosci e
ogni concetto teorico sarà subito collegato al codice reale del progetto
\cite{progetto2026}.

Il progetto sostituisce un'applicazione legacy basata su regole e Q-learning
(archiviata in \file{archive/legacy\_qtable\_rule\_based/} e mai importata
dal nuovo runtime) con un percorso di apprendimento supervisionato completo:
un dataset sintetico etichettato da un sistema a regole, una rete neurale
addestrata a imitarlo e un \emph{safety gate} deterministico che ha l'ultima
parola in ogni decisione.

\section{IA, machine learning e deep learning: cosa significano davvero}\label{sec:ia-ml-deep}

I tre termini che si sentono più spesso — intelligenza artificiale, machine
learning e deep learning — non sono sinonimi: sono tre cerchi concentrici.

\begin{itemize}
  \item \textbf{L'IA è il campo.} È la disciplina che studia come costruire
        sistemi capaci di svolgere compiti che, negli esseri umani,
        richiedono intelligenza: riconoscere un volto, capire una frase,
        decidere se un paziente è a rischio. È il contenitore più ampio.
  \item \textbf{Il machine learning è il metodo statistico.} Dentro il campo
        dell'IA, il ML è l'approccio che non scrive le regole a mano, ma le
        \emph{impara dai dati}. Da programmatore, la differenza è radicale:
        nel software tradizionale scrivi \codice{if valore > soglia: ...}; nel
        ML scrivi un algoritmo che, dato un insieme di esempi, trova da solo
        la funzione di decisione.
  \item \textbf{Il deep learning è una famiglia di modelli.} Dentro il ML, il
        deep learning usa reti neurali con molti strati di trasformazioni.
        È molto potente su dati complessi (immagini, testo, audio), ma non è
        l'unico modo di fare ML e non è sempre la scelta giusta.
\end{itemize}

Un modo sintetico di ricordarlo: \emph{l'IA è il campo, il ML è il metodo
statistico, il deep learning è una famiglia di modelli}
\cite{goodfellow2016deep}.

In questo progetto useremo machine learning \emph{supervisionato} con una
rete neurale piccola (un percettrone multistrato, MLP) scritta da zero in
numpy. Non è deep learning nel senso di migliaia di strati: è una rete a due
strati, sufficiente per un problema con sei ingressi e tre classi. Il punto
del progetto non è la complessità del modello, ma l'intero percorso: dati,
etichette, training, sicurezza e riproducibilità.

\section{Il problema di questo progetto}\label{sec:problema-progetto}

Il dominio è la \textbf{telemedicina}: il monitoraggio remoto dei pazienti.
Un dispositivo misura sei parametri vitali:

\begin{itemize}
  \item pressione sistolica e diastolica (in mmHg);
  \item frequenza cardiaca (in battiti al minuto);
  \item temperatura corporea (in gradi Celsius);
  \item saturazione dell'ossigeno nel sangue (in percentuale);
  \item glicemia (in mg/dL).
\end{itemize}

Il sistema deve classificare ogni combinazione di questi sei valori in una
delle tre classi di rischio: \classe{basso}, \classe{medio} o \classe{alto}.
È un problema di \textbf{classificazione supervisionata}: abbiamo un insieme
di esempi in cui l'output atteso è noto (le etichette) e vogliamo una
funzione che, dato un nuovo input mai visto, produca la classe giusta.

Perché la sicurezza è prioritaria in telemedicina? Perché un errore non è un
bug innocuo: classificare come \classe{basso} un paziente in crisi può avere
conseguenze cliniche. Per questo il progetto adotta una regola architetturale
ferrea: le regole deterministiche di sicurezza hanno \textbf{precedenza
assoluta} sul modello. Un caso che le regole giudicano critico non passa mai
dalla rete neurale e non può mai essere declassato. Il modello impara a
imitare le regole, ma le regole restano il garante finale.

\section{La soluzione in una frase}\label{sec:soluzione-in-una-frase}

La soluzione del progetto si può riassumere in una frase:

\begin{quote}
Un \emph{teacher} rule-based etichetta un dataset sintetico; un MLP in numpy
impara ad approssimare il teacher; a runtime un \emph{safety gate}
deterministico ha precedenza assoluta sul modello.
\end{quote}

In pratica il percorso è questo:

\begin{enumerate}
  \item \textbf{Teacher rule-based}: un sistema a regole (if/elif su soglie
        cliniche) assegna la classe di rischio a ogni combinazione di
        parametri. È la fonte delle etichette (capitolo \ref{cap:teacher}).
  \item \textbf{Dataset sintetico}: il teacher etichetta 70\,000 campioni
        generati artificialmente, congelati in train/val/test
        (capitolo \ref{cap:dati}).
  \item \textbf{MLP numpy}: una rete neurale a due strati impara ad
        approssimare la funzione del teacher (capitoli \ref{cap:reti} e
        \ref{cap:training}).
  \item \textbf{Safety gate}: a runtime, le regole di sicurezza vengono
        applicate prima del modello; il modello interviene solo sui casi che
        le regole non considerano critici (capitolo \ref{cap:sicurezza}).
\end{enumerate}

\begin{esempio}{Il flusso in una riga di codice}
Nel runtime, il metodo \codice{predict} di \codice{SupervisedAgent} applica
prima \codice{analizza} (le regole di sicurezza) e solo dopo, se il caso non
è critico, interroga l'MLP. Il codice reale è in
\file{src/telemedicina\_supervised/agents/supervised\_agent.py}.
\end{esempio}

\section{Mappa del documento}\label{sec:mappa-documento}

Il documento segue il percorso del progetto, dalla teoria al codice.

\begin{description}
  \item[\ref{cap:fondamenti} — Fondamenti di machine learning supervisionato:]
        i concetti base: dati, feature, label, modello come funzione,
        training e test, generalizzazione.
  \item[\ref{cap:dati} — I dati:] le sei feature cliniche, il dataset
        sintetico, lo split congelato e la normalizzazione.
  \item[\ref{cap:teacher} — Il teacher:] il sistema a regole che produce le
        etichette, con i suoi vantaggi e limiti.
  \item[\ref{cap:reti} — Le reti neurali:] come funziona un MLP, dal
        percettrone alla backpropagation.
  \item[\ref{cap:training} — Il training:] l'ottimizzazione dei parametri,
        la loss, la discesa del gradiente e l'early stopping.
  \item[\ref{cap:valutazione} — La valutazione:] come si misura se il modello
        funziona: accuracy, precision, recall, kappa.
  \item[\ref{cap:distillazione} — La distillazione:] il collegamento tra
        teacher e studente, il cuore del progetto.
  \item[\ref{cap:sicurezza} — Sicurezza e incertezza:] il safety gate e la
        soglia di confidenza.
  \item[\ref{cap:riproducibilita} — Riproducibilità:] seed, split congelati e
        hash di provenienza.
  \item[\ref{cap:conclusioni} — Conclusioni:] cosa abbiamo imparato e quali
        sono i limiti dichiarati.
\end{description}

\section{Come leggere il codice}\label{sec:come-leggere-codice}

Il codice del progetto è organizzato in quattro aree principali:

\begin{itemize}
  \item \file{src/telemedicina\_supervised/}: il codice runtime. Contiene le
        regole di sicurezza (\file{safety/safety\_rules.py}), il modello
        (\file{ml/}), l'agente (\file{agents/supervised\_agent.py}), i
        servizi, il database e la configurazione (\file{config.py}).
  \item \file{training/}: il codice offline. Contiene il teacher
        (\file{teacher\_rules.py}), il generatore del dataset sintetico
        (\file{synthetic\_generator.py}), il training
        (\file{train\_model.py}) e le metriche (\file{metrics.py}).
  \item \file{tests/}: la suite di test unitari (unittest).
  \item \file{tools/}: la dashboard di visualizzazione. È l'unico posto in
        cui è ammesso matplotlib, e solo come strumento di sviluppo: il
        runtime non lo importa mai.
\end{itemize}

Per eseguire la suite di test (129 test):

\begin{lstlisting}[caption={Esecuzione della suite di test del progetto}]
PYTHONPATH=src python -m unittest discover -s tests
\end{lstlisting}

Le sei feature, nell'ordine ufficiale del contratto
(\codice{FEATURE\_ORDER} in \file{safety\_rules.py}), sono:

\begin{table}[h]
\centering
\caption{Le sei feature del progetto e le loro unità di misura.}
\label{tab:feature-unita}
\begin{tabular}{ll}
\toprule
\textbf{Feature} & \textbf{Unità di misura} \\
\midrule
\codice{pressione\_sistolica} & mmHg \\
\codice{pressione\_diastolica} & mmHg \\
\codice{frequenza\_cardiaca} & bpm \\
\codice{temperatura} & °C \\
\codice{saturazione\_ossigeno} & \% \\
\codice{glicemia} & mg/dL \\
\bottomrule
\end{tabular}
\end{table}

Nel prossimo capitolo (\ref{cap:fondamenti}) introduciamo i concetti di
machine learning che servono a capire tutto il resto.